\documentclass[11pt]{article}
\usepackage[margin=1in]{geometry}
\usepackage{amsmath,amssymb,amsthm,mathtools}
\usepackage[colorlinks=true,linkcolor=blue,citecolor=blue,urlcolor=blue]{hyperref}
\newtheorem{theorem}{Theorem}
\newtheorem{lemma}[theorem]{Lemma}
\newtheorem{corollary}[theorem]{Corollary}
\theoremstyle{remark}\newtheorem{remark}[theorem]{Remark}
\newcommand{\T}{\mathbb T}
\newcommand{\C}{\mathbb C}
\newcommand{\R}{\mathbb R}
\newcommand{\norm}[1]{\left\lVert#1\right\rVert}
\newcommand{\abs}[1]{\left|#1\right|}
\DeclareMathOperator{\diag}{diag}
\title{Uniform Random Row Subsampling of One-Dimensional\par Multi-Clump Vandermonde Matrices}
\author{Working mathematical draft}
\date{October 7, 2026}
\begin{document}
\maketitle
\begin{abstract}
For a fixed collection of tightly clustered angular nodes, uniform sampling
without replacement preserves its full Vandermonde Gram matrix to relative
error $\rho$. The sufficient row count is
$O(\rho^{-2}\sum_a n_a^2\log(n/\delta))$.
The sampling constant is absolute; the small-clump and large-separation
thresholds may depend on the total number of nodes and the largest clump.
No lower bound on within-clump spacing is needed for the relative Gram result.
Under comparable within-clump spacing, the smallest singular value retains
the scale $(M\Delta)^{n^\star-1}$. We separate the continuous-to-discrete
estimate, clump assembly, and sampling argument into individual lemmas.
\end{abstract}

\section{Setting and main result}
The angular torus is $\T=\R/(2\pi\mathbb Z)$, equipped with
\[
 d_{\T}(x,y)=\min_{p\in\mathbb Z}\abs{x-y+2\pi p}.
\]
All matrices are complex. Singular values are listed in decreasing order,
with $n$ column singular values, including zeros if the rank is smaller than
$n$. The nodes in $\mathcal X$ are distinct points of $\T$.
Write $\mathcal X=\bigsqcup_{a=1}^{A}\mathcal C_a$ for a partition into
nonempty clumps, and put $n_a=\abs{\mathcal C_a}$ and
$S=\sum_{a=1}^{A}n_a^2$. For $\Omega\subseteq\{0,\ldots,M\}$ of cardinality
$m$, define
\[
 A_M=\frac1{\sqrt{M+1}}[e^{ikx_j}]_{k=0,\ldots,M; j=1,\ldots,n},
 \qquad
 A_\Omega=\frac1{\sqrt m}[e^{ikx_j}]_{k\in\Omega; j=1,\ldots,n}.
\]
These definitions are independent of the real representatives chosen for
the angular nodes.

\begin{theorem}[Relative Gram preservation and smallest singular value]
\label{thm:multi-clump-random-sampling}
The absolute constant $C=3072$ has the following property.
For every pair of integers $2\le n^\star\le n$, there exist constants
\[
 0<c_0=c_0(n,n^\star)<1,\qquad C_0=C_0(n,n^\star)\ge n
\]
such that the following holds. Suppose $M\ge C_0$, $\mathcal X\subset\T$
is fixed, $\max_a n_a=n^\star$, and
\begin{align}
 \max_{x,y\in\mathcal C_a}d_\T(x,y)&\le c_0/M &&(1\le a\le A),
 \label{eq:small-clumps}\\
 \min_{x\in\mathcal C_a,\ y\in\mathcal C_b}d_\T(x,y)&\ge C_0/M
 &&(a\ne b).\label{eq:clump-separation}
\end{align}
Choose $\Omega$ uniformly among all $m$-element subsets of
$\{0,\ldots,M\}$. For $0<\rho,\delta<1$ and $1\le m\le M+1$, assume
\begin{equation}
 \label{eq:multi-clump-sampling-rate}
 m\ge \frac{3072}{\rho^2}S\log\frac{2n}{\delta}.
\end{equation}
With probability at least $1-\delta$, one has
\begin{equation}
 \label{eq:relative-gram}
 (1-\rho)A_M^*A_M\preceq A_\Omega^*A_\Omega
 \preceq(1+\rho)A_M^*A_M.
\end{equation}
On the same event, all singular values satisfy
\begin{equation}
 \label{eq:relative-singular-values}
 \sqrt{1-\rho}\,\sigma_j(A_M)\le\sigma_j(A_\Omega)
 \le\sqrt{1+\rho}\,\sigma_j(A_M),\qquad 1\le j\le n.
\end{equation}
Moreover, for every $K\ge1$ there are positive constants
$c_1,C_1$, depending only on $n,n^\star,K$, such that if $\Delta>0$ and
\begin{equation}
 \label{eq:comparable-spacing}
 \Delta\le d_\T(x,y)\le K\Delta
 \quad\text{for distinct }x,y\in\mathcal C_a,
\end{equation}
then, on that same event,
\begin{equation}
 \label{eq:sampled-smallest-singular-value}
 c_1\sqrt{1-\rho}(M\Delta)^{n^\star-1}
 \le\sigma_{\min}(A_\Omega)
 \le C_1\sqrt{1+\rho}(M\Delta)^{n^\star-1}.
\end{equation}
The constants $c_0,C_0$ are chosen before $K,\Delta$ and do not depend on
them. The random event in \eqref{eq:relative-gram} depends only on
$\mathcal X,\rho,m,M$; it works simultaneously for every admissible
$K,\Delta$ and every coefficient vector for this fixed $\mathcal X$.
\end{theorem}

\begin{remark}\label{rem:scope}
The theorem is a fixed-node-set assertion. It does not claim one event
uniformly over all node configurations. If the sufficient rate exceeds
$M+1$, no sampling cardinality satisfies it. The stated implication still
holds; the deterministic full sample also satisfies \eqref{eq:relative-gram}.
We do not insert a full-sampling alternative into the sufficient-rate
hypothesis of Theorem~\ref{thm:multi-clump-random-sampling}.
\end{remark}

\section{The deterministic row estimate}
The proof below uses only finite-dimensional polynomial estimates and
matrix exponentials. In particular, its constants need not remain uniform
in the clump cardinality when choosing geometric thresholds. This permits
us to prove the estimates actually required here without importing an
unformalized continuous exponential-sum or clustered-subspace theorem.

\begin{lemma}[Polynomial evaluation and coefficient control]
\label{lem:polynomial-evaluation}
For $s\ge1$ and $a=(a_0,\ldots,a_{s-1})\in\C^s$, put
$p_a(t)=\sum_{r=0}^{s-1}a_rt^r/r!$. There is $\gamma_s>0$, depending only
on $s$, such that
\[
 \gamma_s\norm a_\infty^2\le\int_0^1\abs{p_a(t)}^2\,dt,
 \qquad
 \sup_{[0,1]}\abs{p_a}^2\le64s^2\int_0^1\abs{p_a(t)}^2\,dt.
\]
\end{lemma}
\begin{proof}
The Chebyshev endpoint derivative estimate, applied after affine rescaling
on a one-sided subinterval of length at least $1/2$, gives
$\abs{P'(t)}\le4(s-1)^2\sup_{[0,1]}\abs P$ for a real polynomial of degree
at most $s-1$. Applying this to the real and imaginary parts gives the
conservative estimate $\sup\abs{p_a'}\le8s^2\sup\abs{p_a}$.
Let $U=\sup\abs{p_a}$. At a maximizer, choose a side interval of length
$1/(16s^2)$ contained in $[0,1]$. The derivative bound implies
$\abs{p_a}\ge U/2$ there, proving the second assertion.
The linear evaluation map $a\mapsto p_a|_{[0,1]}$ is injective: a polynomial
vanishing on an interval is zero. Equivalence of finite-dimensional norms
therefore supplies $H_s>0$ with $\norm a_\infty\le H_s\sup\abs{p_a}$.
Take $\gamma_s=(64s^2H_s^2)^{-1}$.
\end{proof}

\begin{lemma}[Small frequencies, including collisions]
\label{lem:small-frequency-evaluation}
For every $s\ge1$ there exist $\eta_s\in(0,1]$ and $D_s,\alpha_s>0$,
depending only on $s$, such that every
$g(t)=\sum_{j=1}^s b_je^{i\lambda_jt}$ with real
$\abs{\lambda_j}\le\eta_s$ satisfies
\begin{equation}
\label{eq:continuous-pointwise}
 \sup_{[0,1]}\abs g^2\le128s^2\int_0^1\abs g^2,
 \qquad
 \sup_{[0,1]}\abs{g'}\le D_s\sup_{[0,1]}\abs g.
\end{equation}
If $a_r=\sum_jb_j(i\lambda_j)^r$ for $0\le r<s$, then also
$\alpha_s\norm a_\infty^2\le\int_0^1\abs g^2$.
No distinctness or minimum gap is needed.
\end{lemma}
\begin{proof}
Write $\zeta_j=i\lambda_j$ and
$Q_\zeta(z)=\prod_j(z-\zeta_j)=z^s+\sum_{r<s}q_rz^r$.
Let $C_\zeta$ be the companion matrix with superdiagonal entries one,
last row $(-q_0,\ldots,-q_{s-1})$, and all other entries zero.
With $W_{rj}=\zeta_j^r$, one has
$C_\zeta W=W\diag(\zeta_j)$, hence
\[
 e^{tC_\zeta}W=W\diag(e^{t\zeta_j}),\qquad
 g(t)=(e^{tC_\zeta}a)_0.
\]
These identities do not require inversion of $W$ and hold at collisions.
At $\zeta=0$, $C_0$ is the nilpotent shift and
$(e^{tC_0}a)_0=p_a(t)$. Polynomial dependence of the entries of $C_\zeta$
and continuity of the matrix exponential give uniform convergence of the
first row of $e^{tC_\zeta}$ to $(t^r/r!)_{r<s}$ on $0\le t\le1$.
Consequently, for any prescribed $\varepsilon>0$, a sufficiently small
$\eta_s$ ensures
$\sup\abs{g-p_a}\le\varepsilon\norm a_\infty$.
Choose $\varepsilon^2\le\gamma_s/25$, where $\gamma_s$ is from
Lemma~\ref{lem:polynomial-evaluation}. The inequality
$\abs{u+v}^2\le\tfrac54\abs u^2+5\abs v^2$ first gives
\[
 \int\abs{p_a}^2\le\frac{25}{16}\int\abs g^2,
 \qquad \frac{\gamma_s}{2}\norm a_\infty^2\le\int\abs g^2.
\]
Applying the same inequality pointwise, and the polynomial row estimate,
yields $\sup\abs g^2\le128s^2\int\abs g^2$.
Finally, $g'(t)=(e^{tC_\zeta}C_\zeta a)_0$. Compactness of
$\{\norm\zeta_\infty\le1\}\times[0,1]$ bounds the derivative row by a
constant depending only on $s$. The coefficient-energy lower bound and
$\int\abs g^2\le\sup\abs g^2$ prove the derivative assertion.
\end{proof}

\begin{lemma}[Short circular clumps have short real lifts]
\label{lem:short-clump-lift}
If $d_\T(x,y)\le h$ for every pair in a nonempty finite set and
$3h<2\pi$, there are real representatives with ordinary diameter at most
$h$. Their ordinary pairwise distances equal their circular distances if
$h<\pi$.
\end{lemma}
\begin{proof}
Fix an anchor $x_0$ and lift each point into $[x_0-h,x_0+h]$.
Two lifted points have ordinary distance at most $2h$. Their difference
is within $h$ of some $2\pi p$. If $p\ne0$, the triangle inequality gives
$2\pi\le3h$, a contradiction. Thus their ordinary distance is at most
$h$. Distances smaller than $\pi$ realize the circular minimum.
\end{proof}

\begin{lemma}[Single-clump discrete row bound]
\label{lem:single-clump-row}
For every $s\ge1$ there are $0<\eta_s\le1$ and $N_s\ge s$, depending
only on $s$, such that for $M\ge N_s$, nodes in a circular clump of diameter
at most $\eta_s/M$, and $f(k)=\sum_{j=1}^sb_je^{ikx_j}$,
\begin{equation}
 \label{eq:single-clump-row}
 \abs{f(k)}^2\le\frac{512s^2}{M+1}
                 \sum_{\ell=0}^{M}\abs{f(\ell)}^2,
 \qquad 0\le k\le M.
\end{equation}
No minimum within-clump separation is assumed.
\end{lemma}
\begin{proof}
Fix an anchor $x_0$ and nearest real windings $y_j$ with
$\abs{y_j-x_0}\le\eta_s/M$. Put $\lambda_j=M(y_j-x_0)$ and
$g(t)=\sum_jb_je^{i\lambda_jt}$. Integer windings preserve the samples,
and $f(k)=e^{ikx_0}g(k/M)$. Set $I=\int_0^1\abs g^2$ and
$U=\sup\abs g^2$. Lemma~\ref{lem:small-frequency-evaluation} gives
$U\le128s^2I$ and $\sup\abs{g'}\le D_s\sqrt U$.
Thus $T=\abs g^2$ is Lipschitz with constant $2D_sU$, and the Riemann
rectangle comparison, with $Q=\sum_{\ell=0}^{M}T(\ell/M)$, gives
$MI\le Q+2D_sU$. Choose $N_s\ge\max\{s,1,512D_ss^2\}$.
Absorption gives $Q\ge MI/2$, so
$T(k/M)\le256s^2Q/M\le512s^2Q/(M+1)$.
The same calculation, together with the coefficient lower bound, supplies
$\widetilde\alpha_s>0$ with
$\widetilde\alpha_s\norm a_\infty^2\le Q/(M+1)$, uniformly in these
frequencies and bandwidths.
\end{proof}

\begin{lemma}[Polynomial approximation of a short clump]
\label{lem:clump-polynomial-approximation}
For $s\ge1$ and $\epsilon>0$, there exist $r_s,Q_s,L_s>0$, depending
only on $s,\epsilon$, such that if $M\ge Q_s$, a clump with anchor $x_0$
has diameter at most $r_s/M$, and $u$ belongs to its unnormalized
Vandermonde column space, then some polynomial
$p(t)=\sum_{j=0}^{s-1}c_jt^j$ gives a vector
$w_k=e^{ikx_0}p(k/M)$ satisfying
\[
 \norm{u-w}_2\le\epsilon\norm u_2,
 \qquad \sum_j\abs{c_j}\le\frac{L_s}{\sqrt{M+1}}\norm u_2.
\]
\end{lemma}
\begin{proof}
Represent $u$ by its exponential coefficients and use the companion jet
$a$ from Lemma~\ref{lem:small-frequency-evaluation}.
The grid coefficient lower bound gives
$\norm a_\infty\le D_s'\norm u_2/\sqrt{M+1}$ for a constant $D_s'>0$.
The uniform matrix-exponential convergence permits a radius $r_s$ with
first-row error at most $\epsilon/(sD_s')$ in every entry.
Use $c_j=a_j/j!$. Summing the $s$ entrywise errors proves the vector
error bound. Also $\sum_j\abs{c_j}\le s\norm a_\infty$, so $L_s=sD_s'$
works. No inverse frequency gap enters this argument.
\end{proof}

\begin{lemma}[Clump energy comparison]
\label{lem:clump-energy}
For every $2\le n^\star\le n$ there are thresholds $0<c_0<1$,
$C_0\ge n$, depending only on $n,n^\star$, such that, under the geometry
of Theorem~\ref{thm:multi-clump-random-sampling}, every family
$v_a\in\operatorname{range}V_M(\mathcal C_a)$ satisfies
\begin{equation}
 \label{eq:clump-energy-comparison}
 \frac12\sum_a\norm{v_a}_2^2\le\norm{\sum_av_a}_2^2
 \le\frac32\sum_a\norm{v_a}_2^2.
\end{equation}
Here $V_M(\mathcal C_a)$ is unnormalized, with rows $0,\ldots,M$.
We may simultaneously require $M\ge N_s$ for every $1\le s\le n^\star$.
\end{lemma}
\begin{proof}
Put $\epsilon=1/(32n)$ and apply
Lemma~\ref{lem:clump-polynomial-approximation} for every $1\le s\le n^\star$.
Let $L\ge1$ bound all the resulting $L_s$ and $Q_s$. Choose $c_0$ smaller
than every $r_s,\eta_s,1/s,1$, and choose $C_0$ at least
$n,L,64n\pi L^2$, and every $N_s$. Additional positive size-dependent
shortness bounds and bandwidth thresholds can be included in these finite
minima and maxima.

For distinct anchors $x,y$ write $\theta=d_\T(x,y)$ and $z=e^{i(y-x)}$.
Finite Abel summation and monotonicity of $(k/M)^d$ show
\[
 \abs{\sum_{k=0}^{M}z^k(k/M)^d}\le\frac2{\abs{1-z}}
 \le\frac\pi\theta \qquad(d\ge0),
\]
since $\abs{1-z}\ge2\theta/\pi$ for $0\le\theta\le\pi$.
For the polynomial approximants $w,w'$ to $u,v$, this implies
\[
 \abs{\langle w,w'\rangle}
 \le\frac\pi\theta\left(\sum_j\abs{c_j}\right)
                         \left(\sum_j\abs{d_j}\right)
 \le\frac{\pi L^2}{C_0}\norm u_2\norm v_2
 \le\frac1{64n}\norm u_2\norm v_2.
\]
Here $\theta\ge C_0/M$ and $M/(M+1)\le1$ were used. Expanding
$u=w+(u-w)$ and $v=w'+(v-w')$ gives
\[
 \abs{\langle u,v\rangle}
 \le\left(\frac1{64n}+2\epsilon+\epsilon^2\right)
                         \norm u_2\norm v_2
 \le\frac1{2n}\norm u_2\norm v_2.
\]
Expanding the norm square of $\sum_av_a$ and using $2uv\le u^2+v^2$
bounds its difference from $\sum_a\norm{v_a}_2^2$ by
$(A-1)/(2n)\sum_a\norm{v_a}_2^2\le\tfrac12\sum_a\norm{v_a}_2^2$,
since the nonempty clumps satisfy $A\le n$.
\end{proof}

\begin{lemma}[Multi-clump leverage bound]
\label{lem:multi-clump-leverage}
With the thresholds of Lemma~\ref{lem:clump-energy}, for every coefficient
vector $c$ and every $0\le k\le M$,
\begin{equation}
 \label{eq:multi-clump-leverage}
 \abs{(A_Mc)_k}^2\le\frac{1024S}{M+1}\norm{A_Mc}_2^2.
\end{equation}
In particular, the corresponding supremum over $c\ne0$ has this bound.
\end{lemma}
\begin{proof}
Let $f_a$ be the unnormalized signal from coefficients in clump $a$ and
$E_a=\sum_{\ell=0}^{M}\abs{f_a(\ell)}^2$. By
Lemma~\ref{lem:single-clump-row} and Cauchy--Schwarz,
\[
 \abs{\sum_af_a(k)}^2
 \le\left(\sum_a\sqrt{\frac{512n_a^2}{M+1}}\sqrt{E_a}\right)^2
 \le\frac{512S}{M+1}\sum_aE_a
 \le\frac{1024S}{M+1}\sum_{\ell=0}^{M}\abs{\sum_af_a(\ell)}^2.
\]
Normalize both sides by $M+1$. Since $M\ge n$ and the nodes are distinct,
the first $n$ rows form an invertible square Vandermonde matrix; hence
$\norm{A_Mc}_2>0$ for $c\ne0$.
\end{proof}

\section{Random sampling}
\begin{lemma}[Sampling a fixed subspace]
\label{lem:leverage-sampling}
Let $A\in\C^{N\times n}$ have full column rank and satisfy
$\abs{(Ac)_k}^2\le(R/N)\norm{Ac}_2^2$ for all $c,k$, with $R>0$.
For a uniform $m$-element subset $\Omega$ of its rows, put
$A_\Omega=\sqrt{N/m}\,A|_\Omega$. If $0<\rho,\delta<1$,
$1\le m\le N$, and $m\ge3R\rho^{-2}\log(2n/\delta)$, then
\[
 \mathbb P\bigl((1-\rho)A^*A\preceq A_\Omega^*A_\Omega
                       \preceq(1+\rho)A^*A\bigr)\ge1-\delta.
\]
\end{lemma}
\begin{proof}
Whiten by $H=(A^*A)^{-1/2}$ and put $U=AH$, so $U^*U=I$.
Writing $u_k$ for row $k$, the leverage hypothesis gives
$\norm{u_k}_2^2\le R/N$. The population
$X_k=N u_k^*u_k$ satisfies $0\preceq X_k\preceq RI$ and
$N^{-1}\sum_kX_k=I$. The two tails in
\cite[Theorem 2.2]{Tropp} for sampling without replacement imply
\[
 \mathbb P\left(\frac1m\sum_{k\in\Omega}X_k
                         \notin[(1-\rho)I,(1+\rho)I]\right)
 \le2n\exp\left(-\frac{m\rho^2}{3R}\right)\le\delta.
\]
We used the usual Chernoff factor estimates valid for $0<\rho<1$.
Congruence by $H^{-1}$ yields the assertion. Hermitian eigenvalue
monotonicity, followed by taking square roots, also yields the
multiplicative comparison of every singular value.
\end{proof}

\section{Smallest singular value under comparable spacing}
\begin{lemma}[Full-matrix smallest singular value]
\label{lem:full-clump-smallest}
The thresholds in Lemma~\ref{lem:clump-energy} can be chosen independently
of $K$ so that, for each $K\ge1$, positive constants $c_1,C_1$ depending
only on $n,n^\star,K$ satisfy
\[
 c_1(M\Delta)^{n^\star-1}\le\sigma_{\min}(A_M)
 \le C_1(M\Delta)^{n^\star-1}
\]
whenever \eqref{eq:comparable-spacing} holds.
\end{lemma}
\begin{proof}
We first give a constructive single-clump interpolation bound. For a
clump of size $s\ge2$ take short real lifts of diameter at most $c_0/M$,
and assume $M\ge2s$. Set $q=\lfloor M/s\rfloor$ and
$r=M-q(s-1)$. Then $q\ge M/(2s)$ and $r+1\ge M/s$.
For node $x_j$, the trigonometric polynomial
\[
 P_j(x)=\prod_{\ell\ne j}
       \frac{e^{iqx}-e^{iqx_\ell}}{e^{iqx_j}-e^{iqx_\ell}}
\]
is cardinal on the clump and has degree $q(s-1)$. Our shortness choice
ensures $q\abs{x_j-x_\ell}\le\pi$; consequently each denominator has
modulus at least $2q\Delta/\pi$. The coefficient $\ell^1$ norm of $P_j$
is at most $(\pi/(q\Delta))^{s-1}$.
Multiply $P_j$ by
$D_j(x)=(r+1)^{-1}\sum_{h=0}^{r}e^{ih(x-x_j)}$.
The product remains cardinal and has degree at most $M$. Its coefficient
$\ell^2$ norm is at most
$B_s=(\pi/(q\Delta))^{s-1}/\sqrt{r+1}$, by the triangle inequality for
convolution with the coefficient vector of $D_j$.
The matrix of these cardinal coefficient rows is a left inverse of $V_M$,
with operator norm at most $\sqrt sB_s$. Therefore
\[
 \sigma_{\min}(V_M)\ge\frac{\sqrt M}{s}
                \left(\frac{M\Delta}{2\pi s}\right)^{s-1}
 \ge\beta_s\sqrt M(M\Delta)^{s-1},\qquad
 \beta_s=\frac{2^{-s/2}}{s(\sqrt2\pi s)^{s-1}}>0.
\]
This conservative $\beta_s$ is the coefficient used in the formal proof,
which specializes the repository's interpolation-packet construction.
For $s=1$ the column norm is $\sqrt{M+1}$.
Enlarge $C_0$ to cover $2s$ for all $s\le n^\star$, and retain the
size-dependent shortness conditions in Lemma~\ref{lem:clump-energy}.
Write $b=32\pi e$ and take
$d=\min_{1\le s\le n^\star}\min\{1,\beta_sb^{s-1}\}>0$.
Every single clump then satisfies
$\sigma_{\min}(V_M)\ge d\sqrt M(M\Delta/b)^{s-1}$.
All these choices depend only on the permitted cardinalities, not on $K$.
There is a clump of size $n^\star\ge2$, so $0<M\Delta\le c_0<1$.
For $s\le n^\star$, therefore,
$(M\Delta/b)^{s-1}\ge b^{-(n^\star-1)}(M\Delta)^{n^\star-1}$.
Together with $\sqrt{M/(M+1)}\ge1/\sqrt2$ and the lower energy bound
in \eqref{eq:clump-energy-comparison}, this proves the lower estimate;
for example one may take the conservative constant
$c_1=d/(4b^{n^\star-1})$.

For the upper estimate, choose a largest clump and short lifts
$y_j=x_j-\zeta$ with $\abs{y_j}\le K\Delta$. Put $r=n^\star-1$.
The $r$ homogeneous equations
$\sum_jc_jy_j^\ell=0$ for $0\le\ell<r$ have a nonzero solution in
$\C^{n^\star}$; normalize it to $\norm c_2=1$ and extend by zero to
the other clumps. The integral Taylor remainder for the real phase gives
\[
 \abs{e^{it}-\sum_{\ell=0}^{r-1}(it)^\ell/\ell!}
 \le\frac{\abs t^r}{r!}.
\]
Consequently, at every $0\le k\le M$,
$\abs{\sum_jc_je^{iky_j}}\le\sqrt{n^\star}(MK\Delta)^r/r!$.
Averaging its square over the $M+1$ rows proves the upper estimate.
We may use the slightly larger constant $C_1=2\sqrt{n^\star}K^r/r!$;
this also follows from the usual complex exponential Taylor remainder
bound on $\abs t\le1$, the version used in the Lean development.
\end{proof}

\begin{proof}[Proof of Theorem~\ref{thm:multi-clump-random-sampling}]
Choose the thresholds simultaneously as above. Apply
Lemma~\ref{lem:leverage-sampling} with $N=M+1$, $A=A_M$, and
$R=1024S$ from Lemma~\ref{lem:multi-clump-leverage}. Its rate is
exactly \eqref{eq:multi-clump-sampling-rate}. This proves
\eqref{eq:relative-gram} and \eqref{eq:relative-singular-values}.
Combine the latter with Lemma~\ref{lem:full-clump-smallest} to obtain
\eqref{eq:sampled-smallest-singular-value}. The same event works for
every subsequent choice of spacing parameters.
\end{proof}

\begin{corollary}[Sampling order]\label{cor:clump-sampling-order}
Under the same geometry, $S\le nn^\star$. Hence a sufficient sampling
threshold is $O(\rho^{-2}nn^\star\log(n/\delta))$, with an absolute implied
constant; for fixed $\rho$ this is $O_\rho(nn^\star\log(n/\delta))$. The threshold contains no $\Delta^{-1}$.
\end{corollary}
\begin{proof}
Each $n_a^2\le n^\star n_a$ and $\sum_an_a=n$.
\end{proof}

\begin{remark}[Relation to the literature]
The continuous leverage and clustered-subspace estimates of
\cite{EMM,BG,BDGY} motivate the result. The proof above instead uses
polynomial estimates, companion matrices continuous at frequency
collisions, and constructive interpolation packets to establish every
analytic consequence needed here. The matrix Chernoff argument is the
without-replacement inequality of \cite{Tropp}, whose proof is included in
the repository's general library. The complete Lean multi-clump theorem
therefore has no external admitted dependencies. We make no claim that the
full unrestricted original theorems, with their sharper constants, have
been formalized. In particular, we use the absolute sufficient constant
$3072$; its numerical optimization is not part of this result.
\end{remark}

\begin{thebibliography}{9}
\bibitem{BDGY} D.~Batenkov, B.~Diederichs, G.~Goldman, and Y.~Yomdin,
\emph{The spectral properties of Vandermonde matrices with clustered nodes},
Linear Algebra and its Applications 609 (2021), 37--72.
\href{https://arxiv.org/abs/1909.01927}{arXiv:1909.01927}.
\bibitem{EMM} T.~Erd\'elyi, C.~Musco, and C.~Musco,
\emph{Fourier sparse leverage scores and approximate kernel learning},
2020. \href{https://arxiv.org/abs/2006.07340}{arXiv:2006.07340}.
\bibitem{BG} D.~Batenkov and G.~Goldman,
\emph{Single-exponential bounds for the smallest singular value of
Vandermonde matrices in the sub-Rayleigh regime}, 2021.
\href{https://arxiv.org/abs/2107.09326}{arXiv:2107.09326}.
\bibitem{Tropp} J.~A.~Tropp,
\emph{Improved analysis of the subsampled randomized Hadamard transform},
Advances in Adaptive Data Analysis 3 (2011), 115--126.
\href{https://arxiv.org/abs/1011.1595}{arXiv:1011.1595}.
\end{thebibliography}
\end{document}
